%% toc: P2 11 11.8
\section{Sheffer stroke, NOR, and proofs of expressive limitations}
\label{sec:p02-w02-p02-009-sheffer}
\index{Sheffer function}\index{functional completeness!single connective}

The note on the sixteen binary truth functions
(Section~\ref{sec:p01-sixteen-functions}) records that NAND and NOR are each
functionally complete, and Section~\ref{sec:p02-w02-p02-008-sheffer} proves
this for positive arities and shows that among the operations of arity at
most two they are the only ones that generate every operation on their own.
This section settles the same question for a single operation of any arity,
and for a single connective together with the nullary constants, in each of
the two scopes of completeness of \cref{def:p02-w02-p02-008-completeness}.
It also proves non-definability results for single operations that pass
some, but not all, of the tests: an operation can reverse both constant
inputs and still be confined to a proper clone.

\emph{Purpose.}  To decide, from three entries of a truth table and one
symmetry check, whether one connective suffices for every truth function,
and to say exactly what changes when the nullary truth functions are
included.

\emph{Inputs.}  From Section~\ref{sec:p02-w02-p02-008-post}: Boolean
operations, composition, clones and generated clones
(\cref{def:p02-w02-p02-008-operations,thm:p02-w02-p02-008-generated}); the
operation set \(C^+\) of a connective set and the positive-arity convention
(\cref{def:p02-w02-p02-008-arity}); positive-arity and nullary-inclusive
completeness (\cref{def:p02-w02-p02-008-completeness}); the five classes
\(\mathsf T_0,\mathsf T_1,\mathsf S,\mathsf M,\mathsf L\) and their tests
(\cref{def:p02-w02-p02-008-five,thm:p02-w02-p02-008-tests}); the fact that
they are clones (\cref{thm:p02-w02-p02-008-five-clones}); Post's criterion
(\cref{thm:p02-w02-p02-008-criterion}); and the NAND and NOR identities
(\cref{thm:p02-w02-p02-008-nand-nor}).  Notation is that of
\cref{def:p02-w02-p02-008-bits}: truth values are \(0,1\), \(\bar0\) and
\(\bar1\) are the constant tuples, and \(\bar a^{\,c}\) is the complement of
\(\bar a\).

\emph{Conclusions.}  Consequences of reversing the constant inputs
(\cref{thm:p02-w02-p02-009-two-tests}); the characterization of single
generators of every arity (\cref{thm:p02-w02-p02-009-sheffer}) and their
number (\cref{thm:p02-w02-p02-009-count}); the completeness of a single
connective, alone and with \(\bottom\) or \(\topv\), in both scopes
(\cref{thm:p02-w02-p02-009-single}); and ternary instances, positive and
negative (\cref{thm:p02-w02-p02-009-ternary}).  Every proof is in the annex.

\subsection{Single generators}
\label{sec:p02-w02-p02-009-single-generators}

\begin{hbdefinition}{def:p02-w02-p02-009-sheffer-op}{Sheffer operation, Sheffer connective}
An operation \(f\in\mathcal O\) is a \emph{Sheffer operation} iff
\(\langle\{f\}\rangle=\mathcal O\), that is, iff every Boolean operation of
positive arity is obtained from \(f\) and the projections by composition.
A truth-functional connective of positive arity is a \emph{Sheffer
connective} iff its interpreting operation is a Sheffer operation.  A
nullary connective is never called a Sheffer connective; its role is settled
separately in \cref{thm:p02-w02-p02-009-single}(1).
\index{Sheffer operation}
\end{hbdefinition}

By \cref{thm:p02-w02-p02-008-formulas-clone}, \(c\) is a Sheffer connective
iff the one-member connective set \(\{c\}\) is positive-arity functionally
complete.  For an \(n\)-ary \(f\), call \(x\mapsto f(x,\ldots,x)\) its
\emph{diagonal}.  The diagonal is negation exactly when \(f(\bar0)=1\) and
\(f(\bar1)=0\), which is the condition \(f\notin\mathsf T_0\) and
\(f\notin\mathsf T_1\).

\begin{hblemma}{thm:p02-w02-p02-009-two-tests}{proved=annex:p02-w02-p02-009-two-tests}{Reversed constants}
Let \(n\geq1\) and \(f\in\mathcal O_n\).
\begin{enumerate}
\item If \(f(\bar0)=1\) and \(f\in\mathsf M\), then \(f=\kappa^n_1\).  If
  \(f(\bar1)=0\) and \(f\in\mathsf M\), then \(f=\kappa^n_0\).
\item If \(f(\bar0)=1\) and \(f(\bar1)=0\), then \(f\notin\mathsf M\).
\item If \(f(\bar0)=1\), \(f(\bar1)=0\) and \(f\in\mathsf L\), then
  \(f\in\mathsf S\).
\end{enumerate}
\end{hblemma}

\begin{hbtheorem}{thm:p02-w02-p02-009-sheffer}{proved=annex:p02-w02-p02-009-sheffer}{Sheffer operations of every arity}
Let \(n\geq1\) and \(f\in\mathcal O_n\).  Then \(f\) is a Sheffer operation
iff
\begin{enumerate}
\item \(f(\bar0)=1\),
\item \(f(\bar1)=0\), and
\item \(f\) is not self-dual: there is \(\bar a\in\mathbb B^n\) with
  \(f(\bar a^{\,c})=f(\bar a)\).
\end{enumerate}
In particular no unary operation is a Sheffer operation, and an operation
whose diagonal is negation is a Sheffer operation iff it is not self-dual.
\end{hbtheorem}

The three conditions are tested on the truth table with at most
\(2+2^{n-1}\) evaluations: the rows \(\bar0\) and \(\bar1\), and the
\(2^{n-1}\) rows with \(a_1=0\) for self-duality
(\cref{thm:p02-w02-p02-008-tests}(1)).  For \(n=2\) the theorem gives back
\cref{thm:p02-w02-p02-008-sheffer}: the binary operations with
\(f(0,0)=1\) and \(f(1,1)=0\) are \(\mathrm{NAND}\), \(\mathrm{NOR}\),
\((x,y)\mapsto1-x\) and \((x,y)\mapsto1-y\), and the last two are self-dual.

\begin{hbcorollary}{thm:p02-w02-p02-009-count}{proved=annex:p02-w02-p02-009-count}{Number of Sheffer operations}
For \(n\geq1\), the number of Sheffer operations in \(\mathcal O_n\) is
\[
  2^{2^n-2}-2^{2^{n-1}-1}.
\]
It is \(0\) for \(n=1\), \(2\) for \(n=2\), \(56\) for \(n=3\) and
\(16256\) for \(n=4\).  The proportion of Sheffer operations among the
\(2^{2^n}\) members of \(\mathcal O_n\) is
\(\tfrac14-2^{-2^{n-1}-1}\), which increases with \(n\) and tends to
\(\tfrac14\).
\end{hbcorollary}

\subsection{One connective and the nullary constants}
\label{sec:p02-w02-p02-009-scope}

The definition of functional completeness in
Section~\ref{sec:prop-syntax-semantics-04} quantifies over every arity, \(0\)
included (\cref{def:p02-w02-p02-008-completeness}).  The next theorem keeps
the two scopes apart for every connective set with one member, and for one
connective of positive arity together with \(\bottom\) (value \(0\)), with a
nullary verum \(\topv\) (value \(1\)), or with both.

\begin{hbtheorem}{thm:p02-w02-p02-009-single}{proved=annex:p02-w02-p02-009-single}{Single connectives in both scopes}
\begin{enumerate}
\item If \(c\) is a nullary connective, \(\{c\}\) is neither positive-arity
  nor nullary-inclusive functionally complete.
\item Let \(c\) be a connective of arity \(n\geq1\) interpreted by
  \(f\in\mathcal O_n\).  Then \(\{c\}\) is positive-arity functionally
  complete iff \(f\) satisfies the three conditions of
  \cref{thm:p02-w02-p02-009-sheffer}; and \(\{c\}\) is never
  nullary-inclusive functionally complete.
\item Consequently no connective set with exactly one member is
  functionally complete in the sense of
  Section~\ref{sec:prop-syntax-semantics-04}.
\item Let \(c\) and \(f\) be as in part~2.  Each of the following sets has a
  nullary member, so for each of them positive-arity and nullary-inclusive
  functional completeness coincide, and:
  \begin{enumerate}
  \item \(\{c,\bottom\}\) is functionally complete iff \(f(\bar0)=1\) and
    \(f\notin\mathsf L\);
  \item \(\{c,\topv\}\) is functionally complete iff \(f(\bar1)=0\) and
    \(f\notin\mathsf L\);
  \item \(\{c,\bottom,\topv\}\) is functionally complete iff
    \(f\notin\mathsf M\) and \(f\notin\mathsf L\).
  \end{enumerate}
  In case~(a) the closed formula \(c(\bottom,\ldots,\bottom)\) represents
  the constant \(1\); in case~(b) the closed formula
  \(c(\topv,\ldots,\topv)\) represents the constant \(0\).
\item If \(c\) is a Sheffer connective, all three sets of part~4 are
  functionally complete.
\end{enumerate}
\end{hbtheorem}

For \(\mathrm{NAND}\) and \(\mathrm{NOR}\), part~2 restates
\cref{thm:p02-w02-p02-008-nand-nor}(2)--(3): each is positive-arity complete
and neither alone represents a nullary truth function.  Part~4 shows that
adding a constant does not require a Sheffer connective: \(\{\to,\bottom\}\)
is complete because \((0\to0)=1\) and \(\to\) is not affine, while
\(\{\to,\topv\}\) is not, because \((1\to1)=1\), so that every operation
generated lies in \(\mathsf T_1\).

\subsection{Ternary instances and limitations}
\label{sec:p02-w02-p02-009-ternary}

A non-definability claim is proved by exhibiting a clone that contains the
given operations but not the target: by
\cref{thm:p02-w02-p02-008-generated}(5) the generated clone stays inside it,
and by \cref{thm:p02-w02-p02-008-formulas-clone} so does the truth function of
every formula over the connective set.  In the next theorem the operations
are given by their values on \(\bar a=(x,y,z)\in\mathbb B^3\);
\(\operatorname{maj}\) is the majority operation of
Section~\ref{sec:p02-w02-p02-008-clones}.

\begin{hbtheorem}{thm:p02-w02-p02-009-ternary}{proved=annex:p02-w02-p02-009-ternary}{Ternary instances}
\begin{enumerate}
\item The ternary operations \(1\oplus xyz\) and
  \((1\oplus x)(1\oplus y)(1\oplus z)\), the three-place NAND and NOR, are
  Sheffer operations.  Identifying the last two arguments of each gives
  \(\mathrm{NAND}\), respectively \(\mathrm{NOR}\).
\item The operation \(g(x,y,z)=1\oplus\operatorname{maj}(x,y,z)\) satisfies
  \(g(\bar0)=1\) and \(g(\bar1)=0\), and lies in \(\mathsf S\) and in none of
  \(\mathsf T_0,\mathsf T_1,\mathsf M,\mathsf L\).  Hence
  \(\langle\{g\}\rangle\subseteq\mathsf S\): no formula over a connective
  set whose operation set is \(\{g\}\) represents \(\land\), \(\lor\), \(\to\),
  \(\mathrm{NAND}\), \(\mathrm{NOR}\) or a constant operation.
\item The operation \(h(x,y,z)=1\oplus x\oplus y\oplus z\) satisfies
  \(h(\bar0)=1\) and \(h(\bar1)=0\), and lies in \(\mathsf S\) and
  \(\mathsf L\) and in none of \(\mathsf T_0,\mathsf T_1,\mathsf M\).  Hence
  \(\langle\{h\}\rangle\subseteq\mathsf S\cap\mathsf L\): besides the
  operations excluded in part~2, it excludes \(\operatorname{maj}\), which
  is not affine.
\item The operation \(f(x,y,z)=(1-x)\cdot(y\to z)\) is a Sheffer
  operation, and for all \(x,y,z\in\mathbb B\):
  \(f(x,x,x)=1-x\), \(f(x,1-x,x)=0\) and \(f(0,y,z)=(y\to z)\).
\end{enumerate}
\end{hbtheorem}

\begin{hbexample}{def:p02-w02-p02-009-worked}{A ternary Sheffer connective and its definitions}
Let \(c\) be a ternary connective interpreted by
\(f(x,y,z)=(1-x)\cdot(y\to z)\), that is, by \(\neg x\land(y\to z)\).  Its
truth table has the value \(1\) exactly at \((0,0,0)\), \((0,0,1)\) and
\((0,1,1)\).
\begin{enumerate}
\item \emph{Tests.}  \(f(0,0,0)=1\) and \(f(1,1,1)=0\).  At
  \(\bar a=(0,1,0)\), \(f(\bar a)=1\cdot(1\to0)=0\) and
  \(f(\bar a^{\,c})=f(1,0,1)=0\); so \(f\) is not self-dual, and
  \cref{thm:p02-w02-p02-009-sheffer} makes \(c\) a Sheffer connective.
\item \emph{Negation.}  The diagonal is
  \(f(x,x,x)=(1-x)\cdot(x\to x)=1-x\); so \(c(p_1,p_1,p_1)\) represents
  \(\neg p_1\).
\item \emph{A constant.}  The witness \((0,1,0)\) of non-self-duality gives
  \(k(x)=f(x,1-x,x)\), with \(k(0)=f(0,1,0)=0\) and \(k(1)=f(1,0,1)=0\); so
  \(k=\kappa^1_0\), represented by
  \(\theta\defeq c(p_1,c(p_1,p_1,p_1),p_1)\).
\item \emph{Implication.}  \(f(0,y,z)=y\to z\), so
  \(f(k(x),x,y)=(x\to y)\) for all \(x,y\); the formula
  \(\iota\defeq c(\theta,p_1,p_2)\) represents \(p_1\to p_2\).
\item \emph{Conjunction.}  \(x\land y=\neg(x\to\neg y)\), so
  \(p_1\land p_2\) is represented by
  \(c(\iota',\iota',\iota')\) with \(\iota'\) the result of replacing \(p_2\)
  by \(c(p_2,p_2,p_2)\) in \(\iota\).
\end{enumerate}
The identities used in items~2 to~4 are part~4 of
\cref{thm:p02-w02-p02-009-ternary}; its proof also checks the five formulas
on every row.  From
\(\neg\) and \(\land\) every operation of positive arity follows
(\cref{thm:p02-w02-p02-008-neg-and}).  The connective set \(\{c\}\) still
has no closed formula; \(\{c,\bottom\}\) obtains the constant \(1\) as
\(c(\bottom,\bottom,\bottom)\), by \cref{thm:p02-w02-p02-009-single}(4).
\end{hbexample}

\begin{hbwarning}{def:p02-w02-p02-009-limits}{What a single connective does not give}
\begin{enumerate}
\item Reversing both constant rows is necessary and not sufficient: \(g\)
  and \(h\) of \cref{thm:p02-w02-p02-009-ternary} reverse them and are
  confined to \(\mathsf S\).  The self-duality test cannot be omitted.
\item ``NAND alone is functionally complete'' holds in the positive-arity
  scope only.  By \cref{thm:p02-w02-p02-009-single}(3), no one-member
  connective set is complete in the scope of
  Section~\ref{sec:prop-syntax-semantics-04}, which includes the two
  nullary truth functions; a Sheffer connective needs \(\bottom\) or
  \(\topv\) beside it.
\item A Sheffer connective defines every truth function, and the
  definitions can be long: the representation of \(\land\) in
  \cref{def:p02-w02-p02-009-worked} has three copies of \(\iota'\).  The
  results here give no bound on formula size.
\item The characterization concerns operations on the two-element set
  \(\mathbb B\); the lemma behind it uses that the only truth values are
  \(0\) and \(1\), and nothing here applies to more than two truth values.
\end{enumerate}
\end{hbwarning}

\subsection{Relationships and sources}
\label{sec:p02-w02-p02-009-relationships}

\emph{Within this part.}  The single-generator theorem is the one-member case
of Post's criterion (\cref{thm:p02-w02-p02-008-criterion}), with
\cref{thm:p02-w02-p02-009-two-tests} removing the monotone and affine tests.
The binary case is \cref{thm:p02-w02-p02-008-sheffer}; the explicit NAND and
NOR definitions of negation, conjunction and the constants are in
\cref{thm:p02-w02-p02-008-nand-nor}; the general negative and positive bases
are in \cref{thm:p02-w02-p02-008-bases}.  The note on the sixteen binary
functions in Section~\ref{sec:p01-sixteen-functions} states NAND and NOR
completeness without the scope; \cref{thm:p02-w02-p02-009-single}(2)--(3)
supply it.

\emph{Post's terms.}  Post calls a function an \(\alpha\)-, \(\beta\)-,
\(\gamma\)- or \(\delta\)-function according as the one-variable function
obtained by putting one variable for all arguments is the identity, the
constant \(1\), the constant \(0\) or negation \cite[\S\,10,
p.~43]{Post1941}.  So the conditions \(f(\bar0)=1\), \(f(\bar1)=0\) of
\cref{thm:p02-w02-p02-009-sheffer} say that \(f\) is a
\(\delta\)-function.  Post's list of all sets of independent generators of
the complete system begins with the form consisting of a single
non-self-dual \(\delta\)-function \cite[\S\,26, p.~107]{Post1941}.  He
obtains it from the remark that an alternating \(\delta\)-function is
self-dual, so that a sole generator need only be a non-self-dual
\(\delta\)-function \cite[\S\,26, p.~108]{Post1941}; in the present terms
that remark is \cref{thm:p02-w02-p02-009-two-tests}(3), alternating
functions being the affine ones (\cref{thm:p02-w02-p02-008-alternating}).
Post reports that Sheffer introduced the binary functions he calls
rejection and incompatibility, that is, NOR and NAND, and that Nicod used the
latter \cite[\S\,26, p.~117]{Post1941}; and that an earlier study by
Zylinski determined which binary functions are generated by single binary
functions \cite[Introduction, p.~6]{Post1941}.  The count of
\cref{thm:p02-w02-p02-009-count} and the statements about constants in
\cref{thm:p02-w02-p02-009-single} are proved here and are not taken from
Post, who excludes constants from his systems
\cite[Introduction, p.~2, n.~5]{Post1941}.
\index{Post's criterion!single generator}
